\scanpage{103}
\mspage{18}
% archive sheet number: 102
\margin{\deleted{\unsure{SGA}}}
% source lines 1--4
c) Rappelons qu'on conjecture que tt schéma régulier satisfait au th. de pureté cohomologique (SGA 4 $\overline{\underline{\mathrm{XVI}}}$ 3.10 b). Cette propriété est démontrée pour l'instant dans les cas suivants :

% source lines 5--7
1°) $X$ est lisse sur un corps parfait (plus gén., les relations (4.1.1) ou (4.1.2) sont vraies si $Y,Z$ sont lisses sur un corps quelconque), cf. (SGA 4 $\overline{\underline{\mathrm{XVI}}}$ 3.7, 3.9).

% source lines 8--9
2°) $X$ est excellent de caractéristique nulle (SGA 4 XIX 3.2). \deleted{Plus généralement}

% source lines 10--16, interlinear revision of pour aucun
\deleted{Mais pour} Plus généralement, \inserted{on prouve dans loc. cit. que} si $X$ est excellent de car. donnée, \deleted{\unsure{et}} et si on dispose de la résolution des singularités des schémas excellents en cette caractéristique, alors $X$ satisfait au th. de pureté cohomologique. Par contre, \deleted{\unsure{dans ce cas}} pour \deleted{un} \inserted{aucun} schéma régulier d'inégales caractéristiques et de dimension $\geq 2$, le théorème de pureté n'est prouvé à l'heure actuelle.

% interlinear addition below the preceding paragraph
\inserted{d) \unsure{Les|La} condition de 4.1. est stable par localisation \uuline{étale} et par passage à un sous-schéma régulier.}

% source lines 17--20
\uline{Proposition} 4.2. Supposons que le th. de pureté \deleted{soit} cohomologique soit vrai pour $X$, et soit $Y$ un sous-schéma fermé régulier de $X$, purement de codimension $d$. Alors on a, si $\Lambda=\unsure{\mathbb Z/n\mathbb Z|Z/nZ}$, ($n$ premier aux car. rés de $X$)

\[
\underline H^i_Y(\Lambda_X)=
\deleted{\left\{\begin{aligned}\unsure{0}\quad\unsure{i\ne 2d}\\ \mu_n^{\otimes(-d)}\quad\text{si}\end{aligned}\right.}
\unsure{0}\qquad\text{si }i\ne 2d,
\tag*{(4.2.1)}
\]

% source line 22
et on a un isom. canonique

\[
\underline H^{2d}_Y(\Lambda_X)\simeq(\mu_n^{\otimes(-d)})_Y\quad\deleted{\unsure{\text{si}}}
\tag*{(4.2.2)}
\]

(ou encore

\[
\underline H^{2d}_Y(\mu_n^{\otimes d})\simeq\deleted{\unsure{\mathbb Z/n\mathbb Z}}\Lambda_Y\qquad).
\tag*{(4.2.2 bis)}
\]
